% Generated draft: requires mathematical and source review.
\documentclass[11pt]{article}
\usepackage[margin=25mm]{geometry}
\usepackage{fontspec}
\setmainfont{Latin Modern Roman}
\usepackage{amsmath,amssymb,mathtools,unicode-math}
\setmathfont{Latin Modern Math}
\usepackage{enumitem,needspace,xurl,hyperref}
\setlength{\parindent}{0pt}
\setlength{\parskip}{5pt}
\setlength{\emergencystretch}{3em}
\newcommand{\N}{\mathbb{N}}
\newcommand{\Z}{\mathbb{Z}}
\newcommand{\Q}{\mathbb{Q}}
\newcommand{\R}{\mathbb{R}}
\newcommand{\C}{\mathbb{C}}
\newcommand{\F}{\mathbb{F}}
\newcommand{\A}{\mathbb{A}}
\newcommand{\PP}{\mathbb{P}}
\newcommand{\E}{\mathbb{E}}
\newcommand{\eps}{\varepsilon}
\newcommand{\dd}{\,\mathrm{d}}
\begin{document}
\section*{UnsolvedMath ID: 5300023}
\textbf{Source problem code:} catalogue entry
\subsection{Definitions and mathematical statement}
Put $\theta_0=(\sqrt5-1)/2$ and $\lambda=e^{2\pi i\theta_0}$. A standard affine normalization of the quadratic polynomial with a fixed Siegel point of this rotation number is
\[
 P(z)=\lambda z+z^2,
\]
whose critical point is $c=-\lambda/2$. Its Siegel disk $\Delta$ at $0$ is the maximal linearization domain: a conformal map $\phi:\mathbb D\to\Delta$ with $\phi(0)=0$ conjugates $P$ to $w\mapsto\lambda w$. The source concerns the boundary geometry at the critical point.

\textbf{Problems (Lennart Carleson).} Prove the experimentally observed opening-angle behavior of $\partial\Delta$ at $c$, described in the source as ``about $120^\circ$,'' and construct the expected renormalization at that point.

An opening angle requires a geometric definition. If two one-sided boundary tangent rays exist, it might mean their interior angle; other possibilities use limiting sectors or scaling geometry. The source supplies no such definition or tangent-existence hypothesis. It also gives neither an exact angle constant nor a tolerance making ``about $120^\circ$'' a precise inequality. Consequently exact equality to $2\pi/3$ is not silently substituted for the experimental statement.

A local Siegel renormalization usually involves appropriately normalized return maps near the critical point, with return times related to the denominators of the continued-fraction convergents of $\theta_0$ (the golden-mean/Fibonacci pattern). The excerpt does not specify the return domains, coordinate changes, renormalization operator, limit topology, scaling constants, or desired fixed-point assertion. Those are required to state the requested construction precisely; no particular operator is invented.

The supplied summary reports quasisymmetric/self-similar boundary progress through Herman--Yampolsky renormalization and says rigorous fixed-point/scaling work addresses the renormalization aspect, but that an exact $120^\circ$ opening constant was not recorded as proven in its accessible literature. These claims have no precise theorem citations and must not be taken to establish the missing geometric angle definition or exact constant. The record preserves partial-progress status.
\subsection{Short English statement}
For the golden-mean quadratic Siegel disk, justify the observed roughly $120^\circ$ opening at the critical point and construct its local renormalization. The meaning of opening angle, numerical tolerance or exact constant, and the required renormalization data are not specified. Reported scaling progress does not authenticate an exact angle.
\subsection{Sources}
\begin{itemize}
\item Local UnsolvedMath registry, permanent ID 5300023, source code catalogue entry, ``Angle and renormalization at the golden-mean Siegel critical point''; proposer: Lennart Carleson, proposed year 1992. Supplied status partially solved, classification PARTIAL-PROGRESS, updated 2026-08-06.
\item Stony Brook, \emph{Open Problems in Dynamical Systems}, 1992, Carleson 2, \url{https://www.math.stonybrook.edu/open-problems-dynamical-systems}. Reported source-TeX-and-PDF-text extraction and access: 2026-07-29. The record says source-side partial results were preserved but no comprehensive modern resolution check was available. Rights: NEEDS\_REVIEW; public preprint access/source provenance verified but redistribution permission not established.
\item The summary mentions Herman--Yampolsky renormalization, quasisymmetry, self-similarity, and fixed-point geometry without specific papers or theorem statements.
\end{itemize}
\end{document}
